\documentclass[11pt]{article}

\newcommand{\cnum}{Math 245B}
\newcommand{\ced}{Winter 2026}
\newcommand{\ctitle}[4]{\title{\vspace{-0.5in}\cnum, \ced\\Week #1: #2}\author{\vspace{-0.35in}\\#3, #4}}
\usepackage{enumitem}
\usepackage[usenames,dvipsnames,svgnames,table,hyperref]{xcolor}
\usepackage{amsmath, amsfonts}
\usepackage{graphicx} % Required for inserting images
\usepackage{float}
\usepackage{listings}
\usepackage{xcolor}
\usepackage{hyperref}
\usepackage{comment}

\renewcommand*{\theenumi}{\alph{enumi}}
\renewcommand*\labelenumi{(\theenumi)}
\renewcommand*{\theenumii}{\roman{enumii}}
\renewcommand*\labelenumii{\theenumii.}

\definecolor{codegreen}{rgb}{0,0.6,0}
\definecolor{codegray}{rgb}{0.5,0.5,0.5}
\definecolor{codepurple}{rgb}{0.58,0,0.82}
\definecolor{backcolour}{rgb}{0.95,0.95,0.92}
\definecolor{header}{HTML}{205459}
\definecolor{solution}{HTML}{204d52}

\newcommand{\solution}[1]{{{\textcolor{header}{Solution:} \textcolor{solution}{#1}}}}
\setlist[enumerate]{label=(\alph*)}

\lstdefinestyle{mystyle}{
    backgroundcolor=\color{backcolour},
    commentstyle=\color{codegreen},
    keywordstyle=\color{magenta},
    numberstyle=\tiny\color{codegray},
    stringstyle=\color{codepurple},
    basicstyle=\ttfamily\footnotesize,
    breakatwhitespace=false,
    breaklines=true,
    captionpos=b,
    keepspaces=true,
    numbers=left,
    numbersep=5pt,
    showspaces=false,
    showstringspaces=false,
    showtabs=false,
    tabsize=2
}


\begin{document}
\ctitle{\#2}{}{Asher Christian}{006-150-286}
\date{}
\maketitle
\section{Theorem 5.2}
\emph{
    $f$ holomorphic in disc $D$ then  $\exists F$ holomorphic in  $D$, $F'(z) = f(z) \forall z \in D$.
    $F(z) = \int_{\gamma_z}^{}f(w)dw$
}
\section{Corollary 5.3}
\emph{
    $f$ holomorphic in a disc $D$, $\gamma$ piecewise smooth simple closed curve $\mathcal{V} \subset D$ then
    \[
        \int_{\mathcal{V}}^{}f(z)dz = 0
    \] 
}
\\
\emph{
    A domain $\Omega$ is star shaped if $\exists z_0 \in \Omega$ such that $\forall z \in \Omega, [z_0,z] \subset \Omega$. $[z,z_0] =\{tz+(1-t)z_0 : 0 \le t \le 1$.
}
\section{Theorem 5.4}
\emph{
    $f$ holo in a domain $\Omega, D$ is a disc with  $\overline{ D} \subset \Omega$ then
    \[
        \frac{1}{2\pi i} \int_{\substack{\partial D \\ \text{counter clockwise}}}^{}\frac{f(w)}{w-z}dw =
    \begin{cases}
        f(z) & z \in D\\
        0 & z \not\in D\\
        \text{undefined} & z \in \partial D
    \end{cases}
    \] 
}

\section{Theorem 5.5}
\emph{
    Assume $f(z)$ is holomorphic in a domain  $\Omega$, then for all  $z_0 \in \Omega$ $f$ has a power series expansion
    \[
    \sum_{n=0}^{\infty}a_n(z-z_0)^{n} = f(z) \qquad |z-z_0| < r
    \] 
    whenever $r < d(z,\partial \Omega) = \inf_{w \in \partial D} |z_0 - w|$
}
Proof: We know for $|z-z_0| < r \implies f(z) = \frac{1}{2\pi i} \int_{|\zeta - z_0| = r}^{}\frac{f(\zeta)}{\zeta - z_0} d \zeta$.
\\$|z-z_0| < |\zeta - z_0| \forall \zeta \in \partial D$ so
\[
\frac{1}{\zeta - z} = \frac{1}{(\zeta - z_0) - (z-z_0)} = \frac{1}{\zeta - z_0}(1 - \frac{z-z_0}{\zeta - z_0})^{-1}
\] 
\[
= \sum_{n=0}^{\infty}\frac{(z-z_0)^{n}}{(\zeta - z_0)^{n+1}}
\] 
so
\[
f(z) = \frac{1}{2 \pi i}\int_{\zeta - z_0| = 1}^{}f(\zeta) \sum_{n=0}^{\infty}\frac{1}{(\zeta - z_0)^{n+1}}(z-z_0)^{n}d \zeta
\] 
\[
f(z) = \sum_{n=0}^{\infty}\left(\frac{1}{2\pi i}\int_{|\zeta - z_0| = r}^{} \frac{f(\zeta)}{(\zeta - z_0)^{n+1}}d\zeta\right) (z-z_0)^{n}
\] 
because
\[
    \sum_{n=0}^{\infty}\frac{(z-z_0)^{n}}{(\zeta-z_0)^{n+1}} \quad \text{converges uniformly when $|\zeta-z_0| = r$}
\] 

\section{Corollary 5.6}
\emph{
    $f$ holomorphic on a domain $\Omega$ iff  $f$ analytic on  $\Omega$. (has power series expansion at any point)
}

\section{Work exercise 5.7 in notes}

\section{Corollary 5.8}
\emph{
    $f$ holomorphic in a domain $\Omega \implies f \in C^{\infty}$ in $\Omega$ and for all  $n$ and all  $z \in |\Omega|$
     \[
    f^{(n)}(z) = \frac{n!}{2\pi i}\int_{\partial D}^{}\frac{f(\zeta)}{(\zeta - z)^{n+1}}d \zeta \qquad z \in D \subset \overline{ D} \subset \Omega
    \] 
}
Proof by induction. Assume it holds for $n-1$ ,  $n \ge 1$ 
 \[
     \frac{f^{(n-1)}(z+h) - f^{(n-1)}(z)}{h} = \frac{(n-1)!}{2\pi i}\int_{\partial D}^{}f(\zeta)\frac{1}{h}\left(\frac{1}{(\zeta - (z+h))^{n}} - \frac{1}{(\zeta - z)^{n}}\right)d \zeta
\] 
let $A_h(\zeta)$ be the thing inside the large parens\\
Lemma:
\emph{
    $A_h(\zeta) \rightarrow \frac{n}{(\zeta - z)^{n+1}}$ uniformly in $\zeta \in \partial D$ if $z \in D$
}\\
Once  we prove this we can use uniform convergence to conclude the theorem.\\
Proof
\[
\alpha^{n} - \beta^{n} = (\alpha-\beta)(\alpha^{n-1} + \alpha^{n-2}\beta + \dots + \alpha\beta^{n-2} + \beta^{n-1})
\] 
\[
\alpha = \frac{1}{(\zeta - (z+h))} \qquad \beta = \frac{1}{z - \zeta}
\] 
\[
A_h(\zeta) = \frac{1}{h}(\frac{1}{\zeta - (z+h)} - \frac{1}{\zeta - z})\left(\sum_{k=0}^{n-1}\frac{1}{(\zeta-z)^{k}}\frac{1}{(\zeta - (z+h))^{n-1-k}}\right)
\] 

\section{Friday}
\emph{
    $f$ holomorphic on domain $\Omega$ then  $f$ has a power series
    \[
    f(z) = \sum_{n=0}^{\infty}a_n(z_0)(z-z_0)^{n}
    \] 
    converging in $\{z : [z-z_0| < d(z_0, \partial \Omega)\}$
    and 
    \[
    a_n = \frac{1}{2\pi i} \int_{|z-z_0| = r}^{}\frac{f(z)}{(z-z_{0})^{n+1}}dz \qquad 0 < r < d(z_0, \partial \Omega)
    \] 
}
the corollary\\
\emph{
    $f$ holomorphic in  $\Omega$ if and only if  $f$ analytic in $\Omega$
}\\
Corollary 5.8\\
\emph{
    \[
    f^{(n)}(z) = \frac{n!}{2\pi i} \int_{ \partial D}^{} \frac{f(\zeta)}{(\zeta - z)^{n+1}} d \zeta , z \in D \subset \overline{ D} \subset \Omega
    \] 
}

\section{ Corollary 5.9 (Cauchy's Inequalities)}
\emph{
    $f$ holomorphic in  $\Omega$,  $z \in \Omega$ and 
    $r < d(z, \partial \Omega)$ Then
    \[
    |f^{(k)}(z)| \le \frac{k!}{k!}{R^{k}} \sup_{|w-z| = R} |f(w)|
    \] 
}

\section{ Corollary 5.10 (Liouville's Theorem)}
\emph{
    A bounded entire function is constant.
}
\section{Corollary 5.11 Fundamental Theorem of Algebra }
\emph{
    if $p(z)$ is a polynomail of degree $n \ge 1$ then there exists  $z_0 \in \mathbb{C}$ such that $p(z_0) = 0$
}\\
Proof: Suppose not, then $\frac{1}{p(z)}$ is entire, therefore constant since it is bounded since $\lim_{z\to \infty}|p(z)| = +\infty$.

\section{Theorem 5.12 Uniqueness of Analytic Continuation}
\emph{
    Let $f(z)$ be holomorphic in a domain  $\Omega$. And let  $\{a_n\}_{n=1}^{\infty}$ be a sequence of points
    in $\Omega$. 
    \begin{enumerate}
        \item $a_n \rightarrow a \in \Omega$ as $n \rightarrow \infty$ and $a_n \ne a \forall n$
        \item $f(a_n) = 0$ for all  $n$
    \end{enumerate}
    then $f =0 $ in  $\Omega$
}
\\
Corollary if we have $f = g $ on  $a_n$ then  $f = g$ in  $\Omega$ \\
Proof:
\begin{enumerate}
    \item $\exists r > 0$ so that  $f = 0$ on  $\{|z-a| < r\}$. Suppose not, then 
         \[
             f(z)= \sum_{n=n_0}^{\infty}\alpha_n(z-a)^{n}  \qquad \alpha_{n_0} \ne 0
        \] 
        then $f(z)=(z-a)^{n_0}g(z)$ with $g(a) = \alpha_{n_0} \ne 0$ and $g(a_n) = 0 \forall n$ and  $g$ is continuous so  $g(a) = 0$ a contradiction
    \item Connectedness argument.
         $\Omega$ is connected so if  $U \subset \Omega$ and $U$ is open and  $U \ne \emptyset$ and  $\Omega \setminus U$ is open then  $U = \Omega$.\\
         Let  $U = \{z \in \Omega:  \exists r = r(z) > 0, f = 0 \text{ on }\{|w-z|<r_z\} \subset \Omega\}$.
         $a \in U$ by the previous argument and  $U$ is open by definition.
         Take $w \in \overline{ U}$ then thereis a sequence of points $\{w_n\}_{n=1}^{\infty}$ approaching $w$ with  $f(w_n) = 0$so by previous argument  $f(w) = 0$ so $U$ is closed and open in  $\Omega$ and so it is equal to  $\Omega$ since  $\Omega$ is connected.
\end{enumerate}

\section{lemma}
\emph{
    An open set  $\Omega$ is pathwise connected if 
    $\exists z_0 \in \Omega$ so that for all $z \in \Omega$ ther exists a simple curve  $\gamma_z \subset \Omega$ $\gamma_z$ connnecting $z_0$ to $z$
     $\gamma_z(t)$ is continuous.
     $\Omega$ is connected if and only if  $\Omega$ is pathwise connected.
}
Proof: Assume $\Omega$ is path connected but not connected then  $\Omega = U \cup V$  $U, V $ open  $U \cap V = \emptyset, U \ne \emptyset, V \ne \emptyset$.
THen there exists a path and it is connected since it is the continuous image of a connected set ( an interval in  $ \mathbb{R}$). But
$\gamma =(\gamma \cap U) \cup (\gamma \cap V)$ so $\gamma$ is not connected a contradiction.\\
Now assume that $\Omega$ is connected. Choose  $z_0 \in \Omega, z \in \Omega$ let $U = \{w \in \Omega: \exists \gamma_w, \gamma_w(0) = z_{0}, \gamma_w(1)=w, \gamma_w \subset \Omega\}$ $U$ is nonempty since it contains  $z_0$ and $U$ is an open set since open balls are pathwise connected.  $U$ is closed since if  $z \in \overline{ U}$ then there is a path to sometihing nearby  since $\Omega$ is open.
So  $U$ is open and closed in  $\Omega$ and so  $U = \Omega$.

\section{Theorem}
Recall Cauchy's inequalities.  $f$ holomorphic on  $\Omega$ and  $R < d(z, \partial \Omega), z \in \Omega$ 
then
\[
|f^{(n)}(z)| \le \frac{n!}{R^{n}}M(R)
\] 
\[
    M(R) = \sup_{|w-z| \le R} |F(w)|
\] 
I claim that if $ R_1 < R_2$ then $M(R_1) \le M(R_2)$ and equality holds if and only if $f$ is constant
\\
Theorem(Maximum Modulus Thoerem)\\
\emph{
    Let $D = \{z: |z-z_0|< R\}$ $D \subset \overline{ D} \subset \Omega$ and $f $ holo on  $\Omega$
    then $|f(z)| \le \sup_{\partial D}|f(w)|$ for all $z \in D$ and equality holds if and only if  $f$ is constant.
}

\end{document}
